\section{Conclusion et perspectives}\label{sec:conclusion}

Ce projet de semestre ME-401 a permis de concevoir et de réaliser les
premières versions fonctionnelles d'une station météo embarquée pour les
expéditions à la voile dans le cadre du projet Sailowtech.

\subsection{Bilan des travaux}

Les travaux accomplis couvrent l'ensemble de la chaîne, de la conception
jusqu'à la validation préliminaire :

\begin{itemize}
  \item \textbf{Architecture électronique} : conception et routage de deux
    PCB circulaires (étage contrôle et étage capteurs), réception de
    l'ensemble des composants, soudure et assemblage réalisés.

  \item \textbf{Firmware} : développement d'un premier firmware de
    validation (lecture des capteurs, enregistrement SD, communication
    Wi-Fi) et d'un module de vision par ordinateur pour la détection
    de la couverture nuageuse.

  \item \textbf{Design mécanique} : conception d'un boîtier cylindrique
    inspiré de l'abri de Stevenson, réalisation de prototypes par
    impression 3D et découpe laser, intégration de membranes Gore-Tex
    pour tester la ventilation étanche.
\end{itemize}

\subsection{Travaux restants}

Les prochaines étapes pour finaliser et valider la station sont :

\begin{enumerate}
  \item Tester la résistance mécanique du boîtier et son étanchéité
    en conditions simulées (aspersions d'eau, chocs).
  \item Quantifier l'effet de serre du dôme en plexiglas et adapter
    le design si nécessaire.
  \item Évaluer la ventilation passive dans différentes conditions de vent
    et valider que les mesures de température ne sont pas biaisées.
  \item Intégrer et valider l'ensemble de la chaîne de mesure lors d'une
    sortie en mer avec l'équipe Sailowtech.
  \item Améliorer l'algorithme de vision par ordinateur pour les cas
    difficiles (ciel couvert uniforme, contre-jour).
\end{enumerate}

\subsection{Conclusion}

La station conçue répond aux objectifs initiaux : elle est compacte,
modulaire, capable de mesurer la qualité de l'air et les paramètres
météorologiques, et peut communiquer avec les instruments du bateau via
NMEA 2000. L'intégration dans l'écosystème Sailowtech ouvre la voie à
un réseau distribué de voiliers collectant des données environnementales
en haute mer, avec un coût matériel très inférieur aux solutions
professionnelles existantes.
